\vspace{-2mm}
\section{Related Works}
\label{sec:related-works}
    \ecoparagraph{Continuous diffusion distillation.}
    Continuous distillation has two main branches: consistency-like methods~\citep{song2023consistency,song2023improved}, which train students to jump along teacher trajectories, and distribution-matching methods~\citep{yin2024one,yin2024improved,zhou2024score,huang2024flow,gushchin2025inverse,kornilov2025universal, selikhanovych2025one},
    which directly optimize the student distribution, usually with an auxiliary or
    fake model. IDLM follows the distribution-matching view and extends it to
    discrete diffusion.

    \ecoparagraph{Inverse distillation.}
    Inverse distillation directly optimizes a student generator to approximate the
    target distribution, placing it within the family of distribution-matching
    distillation methods. This idea was introduced by IBMD for diffusion bridge models~\citep{gushchin2025inverse}, then
    generalized to continuous-space flow and diffusion
    models in UID framework~\citep{kornilov2025universal}. RSD further extended this view to
    discrete-time diffusion and clarified its relation to DMD: a DMD-like update
    appears after a stop-gradient reduction~\citep{selikhanovych2025one}. IDLM
    adopts the same inverse view, but extends it to discrete space.

    \ecoparagraph{Consistency-style distillation for DLMs.}
    The consistency branch has discrete analogues such as
    SDTT~\citep{deschenaux2024beyond} and Duo-DCD~\citep{sahoo2025the}. However,
    the objective mainly teaches trajectory skipping. In practice, this preserves a
    conditional-independence factorization of the teacher's denoising process, which
    can miss the full joint distribution and favor common modes. Thus, especially in
    the one-step limit, these methods lack guaranteeing recovery of the data
    distribution and remain largely empirical. In contrast, IDLM represents the
    student as a sequence-level mixture, see (\wasyparagraph\ref{sec:loss intuition}), and Theorem~\ref{thm:unique solution} shows that its ideal
    objective is minimized only by the data distribution.
    \nocite{consistency2024rep, abulkhanov2026unifying}

    \ecoparagraph{Distribution matching in discrete space.}
    Discrete distribution-matching methods such as
    Di[M]O~\citep{zhu_dimathttmo_2025}, D-MMD~\citep{hoogeboom2026beyond}, and
    DiDi-Instruct~\citep{zheng2026ultrafast} match teacher and auxiliary signals at
    sampled noising times, corresponding to KL matching of noised marginals,
    \[
        \mathcal L_{\mathrm{DMD}}(\theta)
        =
        \int_0^1 w(t)\,
        D_{\mathrm{KL}}\!\left(p_t^\theta\,\|\,p_t^*\right)\,dt .
    \]
    Here $p_t^\theta$ and $p_t^*$ are the student and target noised marginal
    distributions at time $t$.
    IDLM instead matches full diffusion path measures,
    \[
        \mathcal L_{\mathrm{IDLM}}(\theta)
        =
        D_{\mathrm{KL}}\!\left(\mathbb P^\theta\,\|\,\mathbb P^*\right),
    \]
    so the comparison is trajectory-level rather than marginal-only. A DMD-like
    update appears only after a stop-gradient reduction~\citep{selikhanovych2025one},
    so in masked diffusion, because of the \textsc{subs} parameterization, the
    marginal and path-based losses coincide. In uniform diffusion, this distinction
    matters in practice, as our experiments show.

    \ecoparagraph{Continuous-space language flows.}
    A parallel line changes the state space: FLM/FMLM~\citep{lee2026flow} use
    one-hot flows, while ELF~\citep{hu2026elf} uses embedding-space flows. These
    models can obtain strong GenPPL, but their Entropy can drop relative to other
    baselines. This matters because lower GenPPL may trade off Entropy, for example
    when sweeping the sampling temperature. Moreover, FLM underperforms on recent
    benchmarks~\citep{deschenaux2026language}. Thus, it remains unclear whether a
    continuous or discrete state space is preferable for language modeling.
